\section*{الفصل \ref{ch:g3:temperature-thermometers} --- درجة الحرارة وموازينها}

\begin{solution}{exo:g3:temperature-thermometers:1}
\omterm{def:g3:temperature-thermometers:temperature}{ميزان الحرارة} يعطي العدد نفسه للدفء نفسه، أيًّا كان قارئه ومتى
كان. أما الأيدي فقد خدعتها الأوعية الثلاثة (إذ بدا الماء الفاتر
نفسه دافئًا وباردًا معًا) وخدعها المعدن حين بدا أبرد من الخشب في
\omterm{def:g1:light-and-shadows:shadow}{الظلّ} نفسه.
\end{solution}

\begin{solution}{exo:g3:temperature-thermometers:2}
\omterm{def:g2:ice-water-steam:meltfreeze}{يتجمّد} الماء عند \qty{0}{\celsius} ويغلي عند \qty{100}{\celsius}
--- والعددان \omterm{def:g3:temperature-thermometers:temperature}{بالدرجة السلسيوسية}.
\end{solution}

\begin{solution}{exo:g3:temperature-thermometers:3}
حين يسخن يتمدّد \omterm{def:g3:solids-liquids-gases:liquid}{السائل} المحبوس فيصعد العمود؛ وحين يبرد يتقلّص
فينزل العمود. والأنبوب من الرفاعة بحيث يدفع التمدّد الضئيل
العمودَ صاعدًا دفعًا ظاهرًا.
\end{solution}

\begin{solution}{exo:g3:temperature-thermometers:4}
القراءة قبل الأوان (قبل أن يكفّ العمود عن الحركة)، والقراءة من
زاوية مائلة بدل جعل العين في مستوى قمة العمود.
\end{solution}

\begin{solution}{exo:g3:temperature-thermometers:5}
الماء الغالي: \qty{100}{\celsius}. وغرفة الجلوس الدافئة:
\qty{20}{\celsius}. والثلاجة: \qty{4}{\celsius}. وجسمك:
\qty{37}{\celsius}.
\end{solution}

\begin{solution}{exo:g3:temperature-thermometers:6}
القِدر أسخن، بمقدار $100 - 35 = 65$ درجة --- وهو أسخن بكثير من
أن يُلمس: فهو عند درجة الغليان.
\end{solution}

\begin{solution}{exo:g3:temperature-thermometers:7}
ناقص ثلاث درجات سلسيوسية، وتُكتب \qty{-3}{\celsius}. وهي
\emph{أدفأ} من \qty{-8}{\celsius}: فثماني علامات تحت الصفر أوغل
في \omterm{def:g2:ice-water-steam:meltfreeze}{التجمّد} من ثلاث.
\end{solution}

\begin{solution}{exo:g3:temperature-thermometers:8}
$14 - 6 = 8$: دفئ الصباح بمقدار $8$ درجات.
\end{solution}

\begin{solution}{exo:g3:temperature-thermometers:9}
نزل الليل تحت \qty{0}{\celsius} --- وهو معلم \omterm{def:g2:ice-water-steam:meltfreeze}{تجمّد} الماء.
فالبِرَك المتجمّدة \omterm{def:g3:temperature-thermometers:temperature}{ميزان حرارة} من صنع الطقس.
\end{solution}

\begin{solution}{exo:g3:temperature-thermometers:10}
لا --- فالقيمة \qty{16}{\celsius} في الجانب البارد من اللطيف. لكن جلد
إيما يقارن الممرّ بالحديقة المتجمّدة التي غادرتها توًّا، فيبدو
لها دافئًا. وينبغي أن تثق الأسرة \omterm{def:g3:temperature-thermometers:temperature}{بميزان الحرارة}: فهو يقارن
بعلامات سلّمه الثابتة، لا بالمكان الذي جاء منه توًّا.
\end{solution}

\begin{solution}{exo:g3:temperature-thermometers:11}
يبقى الماء الغالي عند \qty{100}{\celsius} مهما طال غليانه ---
فالمعلم سقف: والتسخين الزائد يصنع بخارًا أكثر، لا ماءً أسخن.
ولهذا بالضبط أمكن اختيار \qty{100}{\celsius} معلمًا موثوقًا من
البداية.
\end{solution}
